\chapter{Thermal Effects on Breakdown and Transition}\label{chapter:breakdown}

Turbulence development in hypersonic boundary layers proceeds through a sequence beginning with linear instability growth and culminating in nonlinear breakdown. Whilst \gls{LST} captures the second Mack mode during early amplification, the subsequent nonlinear dynamics remain difficult to predict with comparable accuracy. Breakdown of coherent structures into turbulence determines the heating, frictional drag, and aerothermal performance of hypersonic vehicles. A characteristic feature is the overshoot in surface quantities such as skin-friction coefficient and Stanton number. Rather than the gradual increase predicted by linear theory, both exceed their fully turbulent values during transition before relaxing downstream \citep{Holden1972}. This behaviour arises from abrupt, localised enhancement of mixing and momentum transport during vortex breakdown, which injects turbulent-like fluctuations into the near-wall region faster than the mean flow equilibrates. This transient amplification of aerodynamic heating poses a significant challenge for \gls{TPS} design and risks localised material failure. Understanding this response remains critical for hypersonic vehicle design across a wide range of flight regimes.

Second-mode instabilities are the dominant mechanism in high-Mach-number boundary layers. Controlled-forcing experiments using harmonic point sources confirmed the second mode as the primary instability and identified subharmonic resonance as a breakdown route \citep{Shiplyuk2003,Bountin2008}. More recent work in the Purdue quiet tunnel identified fundamental resonance as the dominant transition pathway on a Mach~6 flared cone \citep{Chynoweth2014} and established amplitude thresholds for turbulent spot formation \citep{Casper2011,Casper2012}. Low-enthalpy studies, including Mach~21 nitrogen-tunnel experiments \citep{Mironov2000}, revealed four-wave parametric resonance between fundamental and harmonic modes. However, such experiments remain limited by quiet-tunnel availability, perfect-gas conditions, and incomplete access to nonlinear dynamics. High-fidelity \gls{DNS} is necessary to systematically assess thermochemical nonequilibrium effects and wall thermal variations.

Temporal and spatial \gls{DNS} studies have investigated fundamental transition mechanisms in compressible hypersonic boundary layers. For instance, \citet{Mayer2011} computed Mach~3 transition triggered by oblique breakdown, a pathway first identified by \citet{Fasel1993}. Such simulations typically employ either broadband forcing across multiple frequencies or selective excitation of dominant energy modes \citep{AlamSandham2000,Bhaskaran2011}. Whilst the latter strategy does not replicate realistic flight or noisy wind tunnels, it efficiently isolates underlying physics. Fundamental and subharmonic resonance originate from amplification of a two-dimensional instability wave. Once this wave reaches finite amplitude, secondary oblique instabilities emerge, growing either at the same frequency (fundamental resonance) or at half the frequency (subharmonic resonance). The current study examines the transition mechanism through the second excitation method, building on \Cref{chapter:disturbance_growth}. Fundamental \gls{DNS} studies of turbulent boundary-layer behaviour in hypersonic flows were initiated by \citet{Martin2007} and advanced by \citet{Duan2011}, who examined the influence of wall temperature and Mach number on mean flow statistics in cold hypersonic flows. \citet{Franko2013} performed a spatial \gls{DNS} of a Mach~6 \gls{ZPG} boundary layer, examining the skin-friction and heat-transfer overshoot observed in experiments and identifying first-mode oblique breakdown as a plausible transition path. An extensive database of cooled boundary layers up to Mach~14 was later generated by \citet{Zhangetal2018}.

Fewer studies have incorporated real gas effects. \citet{Stemmer2005} simulated hypersonic boundary layer breakdown for reactive and non-reactive flow in a Mach~20 boundary layer. \citet{Marxenetal2013,Marxenetal2014} incorporated fundamental resonance and nonlinear modes in a Mach~10 adiabatic flat-plate flow, concluding that chemical reactions had no significant effect on growth rates. \citet{DiRenzo2021} performed a complete transition study incorporating chemical nonequilibrium and observed good qualitative agreement with previous works. \citet{Passiatoreetal2021} provided a comprehensive analysis of breakdown to turbulence under subharmonic second-mode excitation, identifying that fundamental mode saturation and energy transfer to oblique subharmonic modes lead to $\Lambda$-vortex formation. Whilst finite-rate chemistry had little effect on initial transition stages, the endothermic nature of later stages drained modal energy and delayed transition by approximately 7\%.

Several \gls{DNS} studies have examined real-gas effects on hypersonic boundary layers, though most focused on chemical nonequilibrium. \citet{DiRenzo2021} investigated transition in a Mach~10 hypersonic boundary layer at suborbital enthalpy, capturing the sequence of second-mode amplification, resonance with oblique harmonics, and breakdown, yielding four- to five-fold increases in wall friction and heating. That work also considered finite-rate chemistry in cooled turbulent boundary layers, whilst \citet{Passiatoreetal2021} extended the framework to pseudo-adiabatic walls to isolate chemical nonequilibrium effects. A significant limitation was the neglect of thermal nonequilibrium. More recently, \citet{Williamsetal2025} examined chemistry--turbulence coupling and showed that turbulence drives variations in chemical source terms whilst having minimal impact on turbulence structure. Most studies have been limited to fully thermochemical equilibrium conditions at unrealistically high Mach numbers, leaving understanding of isolated thermal effects incomplete.

The following chapter extends the \gls{LST} analysis to the breakdown stage across \gls{CPG}, \gls{TEQ}, and \gls{TNE} formulations in thermally excited states. All cases are treated as chemically frozen, as the selected conditions do not trigger significant chemical activity \citep{BitterShepherd2015}. At high-enthalpy conditions, deviations from equilibrium become important, and it is necessary to assess how different thermal models affect breakdown onset and progression. Attention is given to the interplay between linear amplification, nonlinear resonance, secondary instability growth, and turbulent spot formation. Emphasis is placed on overshoot behaviour of skin friction and heat transfer, the influence of cold-wall conditions in accelerating breakdown, and the collapse of Mack second modes into fully developed turbulence. The analysis further examines the wall-temperature response of \gls{TNE} flows, where vibrational nonequilibrium modifies breakdown dynamics compared to \gls{CPG} and \gls{TEQ} cases. These results provide detailed characterisation of the nonlinear transition regime under different thermal models and wall conditions, highlighting the consequences for aerothermodynamic design and the challenges in developing predictive tools.

\section{Numerical setup and flow definitions}

This study examines spatially developing, zero-pressure-gradient flat-plate boundary layers at freestream Mach numbers of $M_\infty=6$ and $M_\infty=8$. Flow conditions correspond to the ICAO standard atmosphere at 36~km altitude: $p_\infty=498.5~\text{Pa}$, $T_\infty=239.3~\text{K}$, and five-species air ($X_{N_2}=0.79$, $X_{O_2}=0.21$) yielding a stagnation enthalpy of $H=h_\infty+u_\infty^2/2=5.07~\text{MJ/kg}$. The governing equations are the compressible Navier--Stokes equations in the chemically frozen limit. The computational domain is a rectangular box with streamwise extents of $800\delta^*_0$ (Mach~6) and $1600\delta^*_0$ (Mach~8), wall-normal height of $100\delta^*_0$, and spanwise width of $20\delta^*_0$. The grid is uniform in streamwise and spanwise directions, with wall-normal spacing generated via hyperbolic tangent stretching: $y = L_y \sinh(s_1 \xi)/\sinh(s_1)$ for uniformly distributed $\xi \in [0,1]$ and stretching factor $s_1 = 5.0$. Grid resolutions are given in Tables~\ref{tab:domain_grid_dns}, and \ref{tab:domain_grid_iles}.

\begin{table}[hbt!]
\centering
\caption{Grid resolution and thermal conditions for DNS simulations.}
\begin{tabular}{lccccc}
\toprule
Case & Gas Model & $M_\infty$ & $T_w/T_\infty$ &  ($N_x \times N_y \times N_z$) & Total points \\
\midrule
\multirow{3}{*}{M6Tw18} & CPG & 6 & 18 & $5200 \times 600 \times 275$ & $8.58\times 10^8$ \\
 & TEQ & 6 & 18 & $5200 \times 600 \times 275$ & $8.58\times 10^8$ \\
 & TNE & 6 & 18 & $5200 \times 600 \times 275$ & $8.58\times 10^8$ \\
\bottomrule
\end{tabular}
\label{tab:domain_grid_dns}
\end{table}

\begin{table}[hbt!]
\centering
\caption{Grid resolution and thermal conditions for ILES simulations.}
\begin{tabular}{lccccc}
\toprule
Case & Gas Model & $M_\infty$ & $T_w/T_\infty$ &  ($N_x \times N_y \times N_z$) & Total points \\
\midrule
M6Tw12 & TNE & 6 & 12 & $2600 \times 600 \times 205$ & $3.20\times 10^8$ \\
M6Tw18 & TNE & 6 & 18 & $2300 \times 500 \times 165$ & $1.84\times 10^8$ \\
M8Tw12 & TNE & 8 & 12 & $4600 \times 500 \times 165$ & $3.68\times 10^8$ \\
M8Tw18 & TNE & 8 & 18 & $4600 \times 500 \times 165$ & $3.68\times 10^8$ \\
\bottomrule
\end{tabular}
\label{tab:domain_grid_iles}
\end{table}

Inflow boundary-layer profiles are obtained from a locally self-similar framework solved using Chebyshev collocation, ensuring consistent specification across \gls{CPG}, \gls{TEQ}, and \gls{TNE} conditions. For the \gls{TNE} case, an induced \gls{TFG} similarity solution provides thermally frozen inflow states. Conservative variables are initialised by solving the similarity equations and imposed via inlet pressure extrapolation. The outflow and far-field employ simple extrapolation conditions, whilst spanwise boundaries are periodic. The wall is non-catalytic and isothermal at two temperatures: $T_w=1200~\text{K}$ (cold wall) or $T_w=1800~\text{K}$ (hot wall). Cold walls enhance high-frequency second-mode amplification by thinning the boundary layer and shifting the stability spectrum, thereby accelerating breakdown. Hot walls stabilise acoustic-trapped disturbances and delay transition. This systematic variation examines the combined effects of Mach number and wall temperature on instability amplification, nonlinear breakdown, and turbulence development.

The inflow Reynolds number is set to $Re_{\delta^*_0}=10{,}000$, ensuring comparability across cases and sufficient domain length to capture the full transition sequence. The governing equations are discretised using the WENO-Z scheme for convective terms to suppress non-physical oscillations, combined with 4th order central differences for viscous and metric terms. Time integration employs a 3rd order strong stability preserving Runge--Kutta scheme. Transport properties are evaluated using the Musawi--Sandham viscosity model and vibrational relaxation is calculated using the Millikan--White formulation with the Park correction. Simulations use non-dimensional timesteps of $\Delta t = 0.012$ (Mach~6) and $\Delta t = 0.016$ (Mach~8), maintaining stability whilst resolving high-frequency acoustic waves characteristic of second-mode instability.

\subsection{Disturbance seeding and excitation method}

Transition to turbulence is induced via a suction--blowing forcing strip, following \cite{AlamSandham2000} and extended here to a two-dimensional domain with large-amplitude perturbations. A Gaussian-shaped disturbance is imposed at $x_f$ through a wall mass-flux perturbation:
\begin{equation}
v_w(x,z,t) \;=\; a_0 \, g(x) \, h(z) \, \sum_{i=1}^{n} \sin(2\pi f_i t),
\end{equation}
where the streamwise envelope is
\begin{equation}
g(x) = \exp\left(-\frac{(x-x_f)^2}{2(\delta^*_0)^2}\right),
\end{equation}
and the spanwise modulation is
\begin{equation}
h(z) = 1 + \sin\!\left(\frac{2\pi z}{105\delta^*_0} + \frac{2\pi z}{205\delta^*_0}\right).
\end{equation}
The Gaussian envelope has width $2\sigma^2 = 20(\delta^*_0)^2$ and is centred at $x_f = 25\delta^*_0$ from the domain origin, ensuring sufficient grid resolution for accurate representation. The forcing amplitude is set to $a_0 = 0.1 \, u_\infty$, sufficiently large to trigger transition within the domain whilst remaining representative of realistic disturbance levels. The frequency set is case-dependent: for Mach~6, $\{f_i\} = \{0.15,\,0.16\}$, and for Mach~8, $\{f_i\} = \{0.10,\,0.11\}$. These frequencies are selected using linear stability theory to ensure that second-mode growth reaches its maximum between $x/\delta^*_0 \approx 220$ and $280$, coinciding with the forcing location at $x_f \approx 240\delta^*_0$. This alignment ensures the forcing excites the most naturally amplified disturbances and sustains growth within the intended computational region.

\begin{figure}[t]
    \centering
    \includegraphics[width=9cm]{resources/figures/breakdown/insert_shapefunction_natural.pdf}
    \caption{Spatial structure of the wall mass-flux forcing function.}
    \label{figure:forcing_shapefunction}
\end{figure}

Rather than imposing oblique modes, three-dimensionality is introduced through spanwise modulation, which imposes long- and short-wavelength subharmonics ($\lambda_z=105\delta^*_0$ and $205\delta^*_0$). This approach is justified by the stability analysis in Figure~\ref{figure:stability_plot}, which demonstrates the absence of first-mode instability across all thermal models; consequently, oblique modes arising from first-mode interactions need not be artificially introduced. Instead, $\Lambda$-shaped vortices and streaks emerge naturally, isolating the evolution of two-dimensional second-mode disturbances and their secondary instabilities without prescribing an artificial oblique structure.

\begin{figure}[t]
    \centering
    \includegraphics[width=12cm]{resources/figures/breakdown/result_mach6_stability_diagram.pdf}
    \caption{Linear stability analysis at $\delta_{*}^{0} = 255$ for (a) CPG, (b) TEQ, and (c) TNE models.}
    \label{figure:stability_plot}
\end{figure}

Each simulation spans four consecutive flow-through times ($t_f = L_x/u_\infty$). The first two flow-throughs remove initial transients from inflow forcing and wavepacket adjustment; the solution is restarted and continued for two further flow-throughs whilst accumulating Favre-averaged statistics. This procedure ensures statistics represent the fully developed transitional state, isolated from transient effects. The total integration time is compact due to the high-frequency nature of induced acoustic waves and rapid transition dynamics. The resulting dataset captures the complete sequence from controlled second-mode excitation through nonlinear harmonic and subharmonic interactions to turbulence onset.

\section{Influence of thermal excitations on flow behaviour}

The influence of thermal modelling on hypersonic boundary layer behaviour is investigated by comparing \gls{CPG}, \gls{TEQ}, and \gls{TNE} formulations for the Mach~6 hot-wall configuration. Linear stability analyses reveal that thermal modelling substantially affects second-mode growth rates. \gls{TEQ} exhibits the highest growth rates, exceeding those from \gls{TNE} and \gls{CPG}. This enhanced amplification arises from stronger coupling between vibrationally equilibrated energy modes and acoustic disturbances. The \gls{TNE} formulation introduces vibrational relaxation effects that reduce modal amplification and overall growth rates. The \gls{CPG} model yields the weakest instability growth due to the absence of additional energy channels from thermal excitation. These differences in \gls{LST} provide context for interpreting the nonlinear transition dynamics in the simulations. Transition is characterised through skin-friction coefficient, Stanton number, mean velocity profiles, and velocity fluctuation components. The onset of transition is identified from the sharp rise in skin friction and Stanton number. Reynolds stress components and wall-normal heat flux are examined across the transitional and turbulent regions to assess how the three thermal models differ in turbulence production and to evaluate the mechanisms responsible for the heat-transfer overshoot.


\subsection{Early disturbance amplification}
The streamwise evolution of wall-bounded quantities provides diagnostics for characterising the onset and progression of boundary layer transition. The skin-friction coefficient and Stanton number are defined respectively as
\begin{equation}
    C_f = \cfrac{2\tau_w}{\rho_\infty u_\infty^2},
\end{equation}
\begin{equation}
    St = \cfrac{q_w}{\rho_\infty u_\infty c_{p,tr} (T_{aw}-T_w)},
\end{equation}
where $T_{aw}$ is the adiabatic wall temperature derived from the self-similar solution to ensure accuracy across \gls{TEQ} and \gls{TNE} regimes. For the Stanton number, the translational-rotational temperature $T$ is used for \gls{CPG} and \gls{TNE}, whilst \gls{TEQ} employs a weighted combination of $T$ and $T_v$. The skin-friction coefficient reflects near-wall shear stress and disturbance development, whilst the Stanton number quantifies the nondimensional wall heat flux and thermal boundary-layer response. Both parameters characteristically exhibit overshoot relative to the fully turbulent asymptote during nonlinear breakdown.

For laminar boundary layers, the skin-friction and Stanton number coefficients are derived from self-similarity solutions at each streamwise location to provide a baseline for assessing disturbance growth. Following \citet{White2005}, these are expressed as
\begin{equation}
c_{f,lam} = \sqrt{2} f''(0) \,Re_x^{-1/2}\cfrac{\mu_w \rho_w}{\mu_\infty \rho_\infty},
\end{equation}
where $f$ and $\theta$ are the dimensionless velocity and temperature similarity variables. Both depend on the thermal model; formulations are derived from their respective similarity solutions for \gls{CPG}, \gls{TEQ}, and \gls{TNE}.

For turbulent boundary layers, predictions are benchmarked against the compressible Van Driest II correlation \citep{Sandhametal2014}:
\begin{equation}
    c_{f,turb} = \frac{1}{\lambda} \left[ \frac{0.242 \left( \arcsin\alpha + \arcsin\beta \right)}{\log_{10}(c_f Re_x) - 0.76 \log_{10}(T_w/T_\infty) + 0.41} \right],
\end{equation}
where $\lambda = (\gamma - 1) M_\infty^2 r/2$, $\alpha = (2a^2 - b)/Q$, $\beta = b/Q$, $Q = \sqrt{b^2 + 4a^2}$, $a = \sqrt{\lambda T_\infty/T_w}$, and $b = (1+\lambda)T_\infty/T_w - 1$. The turbulent Stanton number is obtained via the Reynolds analogy:
\begin{equation}
St_{turb} = \frac{c_f}{2Pr_{tr}^{2/3}}.
\end{equation}

\begin{figure}[t]
    \centering
    \fontsize{9}{10}\selectfont
    \includegraphics[width=0.75\textwidth]{resources/figures/breakdown/result_wallquantities_cf_st.pdf}
    \caption[Streamwise evolution of $C_f$ and $St$ for Mach~6, $T_w=1800$ K.]%
    {Streamwise evolution of spatially temporally averaged $C_f$ and $St$ for the Mach~6, $T_w=1800$ K case. Results are shown for \gls{CPG}, \gls{TEQ}, and \gls{TNE} formulations, alongside laminar reference solutions and the Van Driest~II turbulent correlation.}
    \label{figure:thermal_cfst}
\end{figure}

In the present work, laminar and turbulent reference profiles are shown only for the \gls{TNE} case, to avoid excessive clutter in the plots. Comparisons indicate that the laminar and turbulent trends for \gls{CPG} and \gls{TEQ} are very similar to those for \gls{TNE}, although the \gls{TEQ} profiles lie slightly lower. Therefore, the use of \gls{TNE} profiles as representative references is sufficient. 

The skin friction and Stanton number across all cases and comparisons against theoretical values are shown in \cref{figure:thermal_cfst}. In the laminar region, all three models closely follow the laminar reference with minor thermal-model differences. A pronounced second-mode resonance peak emerges, with \gls{TEQ} exhibiting a sharp, well-defined peak. \gls{CPG} shows a similar structure delayed by a  few displacement thickness with slightly higher magnitude, whilst \gls{TNE} displays the most attenuated peak. This progressive delay and attenuation reflect vibrational relaxation extracting energy from the acoustic modes driving second-mode growth. Nonlinear breakdown then accelerates, where both quantities overshoot the turbulent asymptote. \gls{TEQ} exhibits the most dramatic overshoot, followed by \gls{CPG} at a downstream location. \gls{TNE} shows considerably subdued overshoot with a more gradual transition to turbulence, demonstrating that vibrational nonequilibrium moderates breakdown intensity. The transition hierarchy is clear: \gls{TEQ} transitions earliest, \gls{CPG} transitions downstream, and \gls{TNE} transitions latest, delayed by approximately 150 displacement thicknesses relative to \gls{TEQ}. Convergence toward the turbulent asymptote occurs, but the dynamics differ markedly. \gls{TEQ} and \gls{CPG} exhibit sharp transitions with pronounced overshoots followed by rapid relaxation, whilst \gls{TNE} shows a protracted, gentler approach with suppressed overshoot throughout. This sequence demonstrates the profound stabilising influence of thermal nonequilibrium on the nonlinear transition process, contrasting sharply with perfect-gas and thermal equilibrium assumptions.

Selected flow quantities and grid parameters at representative stations throughout the transition process are presented in Table~\ref{table:global_quantities_models}, documenting streamwise evolution across five characteristic phases: near resonance of the two-dimensional second-mode, secondary instability growth, imminent breakdown, end of transition, and the fully turbulent state. Boundary-layer integral parameters including the shape factor $H$ and momentum-thickness Reynolds number $Re_\theta$ characterise velocity profile development, showing progressive steepening as the flow transitions from laminar to turbulent. Thermal quantities including maximum temperature rise $T_{max}/T_e$ and vibrational energy terms $B_q^{TR}$ and $B_q^{V}$ quantify thermal excitation effects; in TNE flows, the latter reveals opposing signs characteristic of translational-rotational and vibrational energy balance. A clear delay in transition stations is evident across the thermochemical models, with TNE consistently requiring greater streamwise distances to progress through successive phases compared to the perfect-gas assumption. The friction Reynolds number $Re_\tau$ increases substantially during breakdown, reflecting intensified near-wall velocity gradients as coherent structures develop and break down. Grid spacings satisfy established DNS resolution requirements for capturing near-wall viscous structures and streamwise coherent features, consistent with previous studies \citep{PoggieEtAl2015}. Similar resolutions were recently achieved using the TENO scheme \citep{Williamsetal2025}.

\newpage
\begin{sidewaystable}[H]
\centering
% \small
\caption{Streamwise evolution of span-averaged and time-averaged dimensionless flow quantities across the transition process for 
CPG, TEQ and TNE simulations.}
\begin{tabular}{lccc ccc ccc ccc ccc}
\toprule
 & \multicolumn{3}{c}{Near resonance of 2D mode} 
 & \multicolumn{3}{c}{Secondary instability} 
 & \multicolumn{3}{c}{Imminent breakdown} 
 & \multicolumn{3}{c}{End of transition} 
 & \multicolumn{3}{c}{Turbulent} \\
\cmidrule(lr){2-4} \cmidrule(lr){5-7} \cmidrule(lr){8-10} \cmidrule(lr){11-13} \cmidrule(lr){14-16}
 & CPG & TEQ & TNE & CPG & TEQ & TNE & CPG & TEQ & TNE & CPG & TEQ & TNE & CPG & TEQ & TNE \\
\midrule
$x/\delta^*_0$ & 256 & 252 & 292 & 312 & 291 & 361 & 400 & 337 & 436 & 505 & 388 & 505 & 789 & 600 & 600 \\
$Re_x \times 10^{-6}$ & 4.16 & 4.29 & 4.53 & 4.72 & 4.68 & 5.22 & 5.60 & 5.14 & 5.97 & 6.65 & 5.65 & 6.66 & 9.49 & 7.77 & 7.61 \\
$Re_\tau$ & 138 & 145 & 189 & 160 & 151 & 166 & 387 & 175 & 240 & 514 & 254 & 514 & 913 & 687 & 824 \\
$Re^*_\delta \times 10^{-4}$ & 1.61 & 1.56 & 1.68 & 1.72 & 1.63 & 1.80 & 1.87 & 1.71 & 1.93 & 2.04 & 1.79 & 2.04 & 2.43 & 2.10 & 2.18 \\
$Re_\theta$ & 1177 & 1132 & 1211 & 1365 & 1200 & 1333 & 1815 & 1294 & 1537 & 2469 & 1437 & 1809 & 4974 & 2123 & 2323 \\
$H$ & 13.62 & 13.22 & 13.96 & 12.20 & 13.11 & 14.00 & 10.64 & 12.59 & 12.69 & 10.33 & 11.73 & 11.03 & 10.24 & 10.04 & 10.71 \\
$M_{a,t,max}$ & 0.10 & 0.08 & 0.09 & 0.11 & 0.08 & 0.07 & 0.16 & 0.08 & 0.09 & 0.15 & 0.09 & 0.14 & 0.14 & 0.11 & 0.15 \\
$B_q^{TR}  \times -10^{-2}$ & 7.17 & 2.93 & 5.11 & 6.23 & 2.88 & 3.78 & 8.98 & 3.00 & 4.72 & 8.24 & 3.41 & 6.99 & 7.37 & 4.21 & 7.23 \\
$B_q^{V} \times 10^{-3}$ & 0.00 & $-5.87$ & 5.58 & 0.00 & $-5.77$ & 5.79 & 0.00 & $-6.01$ & 5.64 & 0.00 & $-6.83$ & 5.69 & 0.00 & $-8.45$ & 5.43 \\
$T_{max}/T_e$ & 3.70 & 3.36 & 3.73 & 3.70 & 3.37 & 3.67 & 3.69 & 3.40 & 3.75 & 3.59 & 3.48 & 3.81 & 3.56 & 3.51 & 3.72 \\
$\delta_{M_t}/\delta^{*}_{0}$ & 0.88 & 0.90 & 0.90 & 0.91 & 0.90 & 0.90 & 0.47 & 0.89 & 0.86 & 0.37 & 0.87 & 0.67 & 0.48 & 0.66 & 0.33 \\
\midrule
$\Delta x^+$ & 7.72 & 5.49 & 6.56 & 7.99 & 5.47 & 5.31 & 12.14 & 5.88 & 6.72 & 11.47 & 6.67 & 10.05 & 10.66 & 8.35 & 10.97 \\
$\Delta y_w^+$ & 0.68 & 0.48 & 0.58 & 0.70 & 0.48 & 0.47 & 1.07 & 0.52 & 0.59 & 1.01 & 0.59 & 0.88 & 0.94 & 0.73 & 0.96 \\
$\Delta y^+_{99}$ & 1.54 & 1.52 & 1.97 & 1.74 & 1.58 & 1.72 & 4.00 & 1.81 & 2.46 & 5.22 & 2.59 & 5.20 & 9.15 & 6.89 & 8.27 \\
$\Delta z^+$ & 3.59 & 2.55 & 3.05 & 3.71 & 2.54 & 2.47 & 5.63 & 2.73 & 3.12 & 5.33 & 3.10 & 4.67 & 4.95 & 3.88 & 5.09 \\
$L_z^+$ & 1004 & 714 & 853 & 1038 & 711 & 690 & 1578 & 764 & 873 & 1491 & 867 & 1306 & 1385 & 1085 & 1426 \\
\bottomrule
\end{tabular}
\label{table:global_quantities_models}
\end{sidewaystable}
\newpage

\begin{figure}[H]
    \centering
    \includegraphics[width=22.3cm,angle=-90]{resources/figures/breakdown/result_dns_thermal_stanton_contour.pdf}
    \caption[Time-averaged wall Stanton number ($St$) during boundary layer transition for (a) CPG, (b) TEQ, and (c) TNE models.]%
    {Time-averaged wall Stanton number ($St$) during boundary layer transition for (a) \gls{CPG}, (b) \gls{TEQ}, and (c) \gls{TNE} models.}
    \label{figure:thermal_spanwise_st}
\end{figure}

The spanwise organisation of heat transfer during transition, as shown in Figure~\ref{figure:thermal_spanwise_st}, reveals noticeable differences across the three thermochemical models. The figure displays contours of Stanton number ($St$) as a function of streamwise position and spanwise coordinate, spanning the transitional region between $x/\delta^*_0 = 200$ and $550$. Streak-like structures emerge as transition progresses, redistributing heat flux into alternating high- and low-intensity bands that mark secondary instability growth and rapid breakdown, consistent with observations in prior transition studies \citep{Franko2013}. The spatial organisation and intensity of these streaks differ markedly across models. \gls{TEQ} develops well-defined streaks early in transition, whilst \gls{CPG} exhibits comparable streak organisation with marginally higher peak intensities during breakdown. Both models maintain organised patterns well into the turbulent region. \gls{TNE} exhibits markedly delayed and attenuated streak development, with high-temperature regions emerging considerably further downstream and remaining less intense than in either \gls{CPG} or \gls{TEQ}. By the end of transition, streak organisation weakens across all models and wall heat flux becomes more uniformly distributed, consistent with fully developed turbulence. This progression from well-organised streaks in \gls{TEQ} and \gls{CPG} to attenuated, delayed streaks in \gls{TNE} reflects the stabilising influence of vibrational nonequilibrium on coherent-structure formation and wall heat transfer.

\subsection{Transition structures}
The initial phase of transition is characterised by the amplification of two-dimensional disturbances dominated by the Mack second mode. This behaviour is evident in the modest elevation of skin-friction and Stanton number upstream of breakdown, marking the resonance of the primary instability wavepacket. Linear stability theory predictions presented in Chapter~\ref{chapter:disturbance_growth} demonstrate that the second mode exhibits the largest amplification rates at the frequencies present in this flow. Figure~\ref{figure:lstdns_amplitude} presents linear amplitude growth curves for the representative cases studied herein. The forced frequencies $f = 0.15$ and $f = 0.16$ exhibit peak amplitude growth at distinct streamwise locations, with the observed resonance peaks in the wall-bounded quantities aligning closely with these linear growth regions. This alignment confirms the transition from linear to weakly nonlinear dynamics as energy from the primary modes redistributes into secondary instabilities.

Instantaneous snapshots reveal the coherent structures driving transition. Firstly, the Q-criterion isosurfaces coloured by total density $\rho / \rho_e$, with pressure contours $p/(\rho_e \ u_e^2 )$ in the background, show how instability modes develop and interact across the three thermal models, as shown in Figure~\ref{figure:transition_isosurfaces}. Initially, two-dimensional Mack modes develop as elongated streamwise-aligned structures. 

\newpage
\begin{figure}[H]
    \centering
    \includegraphics[width=10.6cm,angle=0]{resources/figures/test/full_isosurfaces.pdf}
    \caption[Q-criterion isosurfaces coloured by the total density $\rho / \rho_{e}$ during transition for CPG, TEQ, and TNE models.]%
    {Q-criterion isosurfaces coloured by density during transition for (a) \gls{CPG}, (b) \gls{TEQ}, and (c) \gls{TNE} models. Pressure field shown in background plane. Isosurfaces reveal two-dimensional Mack mode development in the initial domain region and subsequent three-dimensional modal interactions driving breakdown.}
    \label{figure:transition_isosurfaces}
\end{figure}
\newpage


% Instantaneous snapshots reveal the coherent structures driving transition. Q-criterion isosurfaces coloured by total density, with pressure contours in the background, show how instability modes develop and interact across the three thermal models, as shown in Figure~\ref{figure:transition_isosurfaces}. Initially, two-dimensional Mack modes develop as elongated streamwise-aligned structures. These grow in amplitude downstream as linear disturbances. The interaction between the two-dimensional Mack mode and three-dimensional secondary modes produces the resonance rise observed in skin-friction and Stanton-number distributions. This resonance represents the transition from linear to weakly nonlinear dynamics, transferring energy from the primary mode into secondary instabilities. The resulting secondary structures drive the streak formation and localised heat-transfer intensification seen in the spanwise distributions, with the strength and organisation of these streaks determining the character of the heat-flux patterns during breakdown. The three-dimensional breakdown of the initially quasi-two-dimensional laminar flow structure is driven by the sustained periodic forcing at $f = 0.15$ and $f = 0.16$, which continues to amplify the primary modes whilst simultaneously exciting oblique secondary waves. The spanwise organisation of wall Stanton number presented in Figure~\ref{figure:thermal_spanwise_st} provides direct evidence of this process: the \gls{TEQ} case develops sharp, well-defined streaks commencing near $x/\delta_0 \approx 300$, whilst the \gls{TNE} case exhibits significantly delayed and attenuated streak development with high-intensity regions appearing substantially further downstream.
The resulting resonance rise observed in skin-friction and Stanton-number distributions represents the transition from linear to weakly nonlinear dynamics, transferring energy from the primary mode into secondary instabilities. The secondary structures drive the streak formation \citep{Klebanoff} and localised heat-transfer intensification seen in the spanwise distributions, with the strength and organisation of these streaks determining the character of the heat-flux patterns during breakdown.

% These grow in amplitude downstream as linear disturbances.  The interaction between the two-dimensional Mack mode and three-dimensional secondary modes produces the resonance rise observed in skin-friction and Stanton-number distributions. This resonance represents the transition from linear to weakly nonlinear dynamics, transferring energy from the primary mode into secondary instabilities. The resulting secondary structures drive the streak formation and localised heat-transfer intensification seen in the spanwise distributions, with the strength and organisation of these streaks determining the character of the heat-flux patterns during breakdown.

% The three-dimensional structure of the boundary layer during transition is further illustrated through instantaneous streamwise velocity snapshots in the $xz$ plane at $y^+ = 1$, presented in Figure~\ref{figure:velocity_snapshot}. The velocity field directly reveals the development of low-speed and high-speed streaks that characterise secondary instability growth across the three thermal models. The \gls{TEQ} case exhibits pronounced, coherent streak formation beginning at comparatively early streamwise locations, with alternating regions of reduced and elevated streamwise velocity apparent by approximately $x/\delta_0 \approx 300$. This early and vigorous streak development is consistent with the rapid secondary instability growth driven by sustained forcing at $f = 0.15$ and $f = 0.16$. The \gls{CPG} case shows similar streak structures but with emergence occurring further downstream, reflecting the slightly delayed transition onset compared to \gls{TEQ}. In contrast, the \gls{TNE} case exhibits considerably delayed and attenuated streak development, with clear velocity fluctuations remaining poorly organised through a substantial portion of the domain. The velocity streaks in \gls{TNE} emerge at substantially more advanced streamwise locations and maintain weaker spatial coherence, demonstrating the stabilising effect of vibrational nonequilibrium on secondary instability growth and coherent structure formation.

The three-dimensional structure of the boundary layer during transition is further illustrated through instantaneous streamwise Q-criterion isosurfaces presented in Figure~\ref{figure:transition_isosurfaces}. The Q-criterion directly reveals the development of low-speed and high-speed streaks characteristic of secondary instability growth across the three thermal models. H-type transition, the most commonly observed phenomenon in hypersonic boundary layers \citep{DeTullio2013}, is driven by the destabilisation of two-dimensional Mack modes \citep{Fedorov2011} and their interaction with oblique second modes \citep{Franko2013, Passiatore2021}. The \gls{TEQ} case exhibits pronounced, coherent streak formation beginning at comparatively early streamwise locations, with alternating regions of reduced and elevated streamwise velocity apparent at early downstream positions. This vigorous streak development is consistent with rapid secondary instability growth driven by sustained forcing at $f = 0.15$ and $f = 0.16$. The \gls{CPG} case shows similar streak structures but with emergence occurring further downstream, reflecting slightly delayed transition onset. In contrast, the \gls{TNE} case exhibits considerably delayed and attenuated streak development, with velocity fluctuations remaining poorly organised through a substantial portion of the domain. Whilst \gls{TEQ} and \gls{CPG} exhibit enhanced mode interactions leading to direct breakdown to turbulence, the \gls{TNE} case displays characteristic $\Lambda$-vortex structures consistent with observations in adiabatic hypersonic flows \citep{Passiatore2021}, demonstrating the stabilising effect of vibrational nonequilibrium on secondary instability growth and coherent structure formation.

% The velocity field directly reveals the development of low-speed and high-speed streaks characteristic of secondary instability growth across the three thermal models. H-type transition, the most commonly observed phenomenon in hypersonic boundary layers \citep{DeTullio2013}, is driven by the destabilisation of two-dimensional Mack modes \citep{Fedorov2011} and their interaction with oblique second modes \citep{Franko2013, Passiatore2021}. The \gls{TEQ} case exhibits pronounced, coherent streak formation beginning at comparatively early streamwise locations, with alternating regions of reduced and elevated streamwise velocity apparent at early downstream positions. This vigorous streak development is consistent with rapid secondary instability growth driven by sustained forcing at $f = 0.15$ and $f = 0.16$. The \gls{CPG} case shows similar streak structures but with emergence occurring further downstream, reflecting slightly delayed transition onset. In contrast, the \gls{TNE} case exhibits considerably delayed and attenuated streak development, with velocity fluctuations remaining poorly organised through a substantial portion of the domain. Whilst \gls{TEQ} and \gls{CPG} exhibit enhanced mode interactions leading to direct breakdown to turbulence, the \gls{TNE} case displays characteristic $\Lambda$-vortex structures consistent with observations in adiabatic hypersonic flows \citep{Passiatore2021}, demonstrating the stabilising effect of vibrational nonequilibrium on secondary instability growth and coherent structure formation.

The three-dimensional structure of the boundary layer during transition is further illustrated through instantaneous streamwise velocity snapshots in the $x-z$ plane at $y^+ = 1$, presented in Figure~\ref{figure:velocity_snapshot}.

Comparison across the three thermal models reveals important differences in the spatial evolution of these structures. For \gls{CPG} and \gls{TEQ}, Mack modes develop and reach significant amplitude in the earlier part of the domain, consistent with their earlier transition onset. The interaction between two-dimensional and three-dimensional modes occurs at comparable streamwise locations, driving the observed overshoot in wall quantities. In contrast, the \gls{TNE} case exhibits delayed mode development, with Mack modes emerging further downstream and reaching large amplitude at a more advanced streamwise location. Despite these spatial delays, the underlying mechanisms of mode interaction and secondary instability growth remain consistent across all cases, indicating that vibrational nonequilibrium modulates the spatial progression of transition without fundamentally altering the breakdown pathway. 

% The differences in modal amplitude and spatial extent reflect the moderated growth rates and stabilising effects of vibrational relaxation discussed in the linear stability analysis, directly linking the instantaneous flow structures to the mean wall-bounded quantities observed in Chapter~\ref{chapter:disturbance_growth}.




% The three-dimensional structure of the boundary layer during transition is further illustrated through instantaneous streamwise velocity snapshots in the $xz$ plane at $y^+ = 1$, presented in Figure~\ref{figure:velocity_snapshot}. The velocity field directly reveals the development of low-speed and high-speed streaks that characterise secondary instability growth across the three thermal models. The \gls{TEQ} case exhibits pronounced, coherent streak formation beginning at comparatively early streamwise locations, with alternating regions of reduced and elevated streamwise velocity apparent by approximately $x/\delta_0 \approx 300$. This early and vigorous streak development is consistent with the rapid secondary instability growth driven by sustained forcing at $f = 0.15$ and $f = 0.16$. The \gls{CPG} case shows similar streak structures but with emergence occurring further downstream, reflecting the slightly delayed transition onset compared to \gls{TEQ}. In contrast, the \gls{TNE} case exhibits considerably delayed and attenuated streak development, with clear velocity fluctuations remaining poorly organised through a substantial portion of the domain. The velocity streaks in \gls{TNE} emerge at substantially more advanced streamwise locations and maintain weaker spatial coherence, demonstrating the stabilising effect of vibrational nonequilibrium on secondary instability growth and coherent structure formation.

% The pronounced differences in transition progression across thermal models reflect fundamental distinctions in how energy redistributes during breakdown. The \gls{TEQ} case undergoes rapid and intense nonlinear breakdown, with the pronounced overshoot in both $C_f$ and $St$ observed approximately $x/\delta_0 \approx 350$--$400$ reflecting swift energy transfer from organised two-dimensional structures to small-scale turbulent features. The \gls{CPG} case exhibits intermediate behaviour with delayed but substantial breakdown dynamics. The \gls{TNE} case demonstrates the most stabilised response, exhibiting delayed and attenuated secondary instability growth with a more gradual transition to turbulence. In the \gls{TNE} case, the vibrational energy mode acts as a thermal energy sink, withdrawing energy from the acoustic modes driving transition. This redistribution into vibrational degrees of freedom moderates the nonlinear breakdown process, resulting in lower transient heating rates and suppressed overshoot intensity. These observations demonstrate that thermal nonequilibrium substantially alters not only the streamwise location of transition but fundamentally modifies the character of the transition process itself, with vibrational relaxation serving as a stabilising mechanism that suppresses the intensity of coherent structure formation and delays the onset of nonlinear breakdown.

\newpage
\begin{figure}[H]
    \centering
    \includegraphics[width=22.3cm,angle=-90]{resources/figures/breakdown/result_contour_yplus1_uvelocity.pdf}
    \caption[Instantaneous streamwise velocity ($u/u_e$) during boundary layer transition for CPG, TEQ, and TNE models.]%
    {Instantaneous streamwise velocity ($u/u_e$) during boundary layer transition for (a) \gls{CPG}, (b) \gls{TEQ}, and (c) \gls{TNE} models.}
    \label{figure:velocity_snapshot}
\end{figure}
\newpage




% \begin{figure}[H]
%     \centering
%     \includegraphics[width=10.3cm,angle=0]{resources/figures/test/full_isosurfaces.pdf}
%     \caption{Linear stability analysis at $\delta_{*}^{0} = 255$ for (a) CPG, (b) TEQ, and (c) TNE models.}
%     \label{figure:thermal_spanwise_st}
% \end{figure}

% \begin{figure}[H]
%     \centering
%     \includegraphics[width=10.3cm,angle=0]{resources/figures/test/test_plot.png}
%     \caption{Linear stability analysis at $\delta_{*}^{0} = 255$ for (a) CPG, (b) TEQ, and (c) TNE models.}
%     \label{figure:thermal_spanwise_st}
% \end{figure}


% \newpage

% \begin{figure}[H]
%     \centering
%     \includegraphics[width=22.3cm,angle=-90]{resources/figures/breakdown/result_contour_yplus1_uvelocity.pdf}
%     \caption[Instantaneous ($u/u_e$) during boundary layer transition for (a) CPG, (b) TEQ, and (c) TNE models.]%
%     {Instantaneous ($u/u_e$)  during boundary layer transition for (a) \gls{CPG}, (b) \gls{TEQ}, and (c) \gls{TNE} models.}
%     \label{figure:thermal_spanwise_st}
% \end{figure}

% \newpage



% \subsection{Forcing interaction}

\subsubsection{Frequency content}

The wall pressure field is analysed through complementary methods to understand the modal mechanisms and spatial organisation during transition: fast Fourier transform (FFT) applied to span-averaged pressure slices to identify dominant frequencies, spatial distribution of pressure RMS to characterise the streamwise evolution of fluctuation intensity, and proper orthogonal decomposition (POD) applied to two-dimensional spanwise-varying pressure data to decompose coherent structures.

Two Mack second-mode frequencies are excited in the forcing, selected from linear stability theory predictions: $f_1 = 0.15$ and $f_2 = 0.16$ (in nondimensional form). The frequency content of wall pressure fluctuations, shown in Figure~\ref{figure:fft_pressure}, reveals the response of each model to these dual-frequency excitation. In the laminar region, pressure fluctuations remain minimal, with power concentrated at low frequencies. As the domain progresses and the second-mode resonance develops, distinct peaks emerge in the power spectral density at frequencies corresponding to the forced modes. Examining the relative amplitudes of the two forcing frequencies provides insight into which mode is preferentially amplified by the boundary layer stability characteristics. \gls{TEQ} exhibits stronger amplification of both frequencies relative to \gls{CPG}, with the spectral peaks more clearly defined and persisting to higher streamwise locations. The dominance of one frequency over the other varies across models and streamwise locations, reflecting the complex interaction between the modal response and the evolving boundary layer structure during the transition process.

\begin{figure}[t]
    \centering
    \includegraphics[width=13.3cm]{resources/figures/breakdown/result_dns_thermal_fft.pdf}
    \caption[Fourier transform of pressure wall data.]%
    {Fourier transform of pressure wall data.}
    \label{figure:thermal_fft}
\end{figure}


The spatial distribution of pressure RMS along the streamwise direction characterises how the intensity of wall pressure fluctuations evolves during transition. Figure~\ref{figure:fft_pressure} (lower panel) shows $p_{rms}$ distributions for \gls{CPG} and \gls{TEQ}. In the laminar region ($x/\delta^*_0 < 200$), both models exhibit low and comparable pressure RMS levels, consistent with the minimal disturbance amplitudes in this regime. As the second-mode resonance develops around $x/\delta^*_0 \approx 250$--$350$, pressure RMS begins to rise, with \gls{TEQ} showing elevated levels relative to \gls{CPG}. The peak in pressure RMS occurs during the nonlinear breakdown phase ($x/\delta^*_0 \approx 350$--$450$), where secondary instabilities reach maximum amplitude. \gls{TEQ} sustains higher $p_{rms}$ throughout this region, consistent with its stronger modal amplification observed in the FFT. Following the peak, pressure RMS decays gradually in the turbulent region as energy cascades toward smaller scales and distributes across a broader frequency spectrum. The streamwise location of peak pressure RMS differs between models, occurring earlier in \gls{TEQ} and displaced downstream in \gls{CPG}, directly reflecting the transition delay observed in wall-bounded quantities. These spatial distributions confirm that thermal equilibrium enhances the intensity of pressure fluctuations throughout the transition process, consistent with stronger coherent structure formation.

\subsubsection{Modal decomposition}

To understand the spanwise variation of pressure fluctuations and identify the dominant coherent structures, proper orthogonal decomposition is applied to wall pressure data on a two-dimensional x-z slice (streamwise and spanwise directions) at fixed times. The energy distribution across the first 40 POD modes, shown in Figure~\ref{figure:pod_energy}, reveals significant differences between \gls{CPG} and \gls{TEQ}. The first POD mode contains approximately 11.5\% of the total kinetic energy in \gls{TEQ} compared to 11\% in \gls{CPG}, indicating that the dominant coherent structure is comparable in both cases but carries slightly higher energy in the thermally equilibrated case. However, the cumulative energy distribution diverges beyond the first few modes: \gls{TEQ} reaches 50\% cumulative energy by mode 8--9, whilst \gls{CPG} requires modes up to approximately 15--18 to achieve the same threshold. This difference indicates that \gls{TEQ} organises pressure fluctuations into fewer, more energetic modes, reflecting the strong coherent streak structures observed in the spanwise heat-transfer distributions. In contrast, \gls{CPG} distributes energy across a broader range of modes, suggesting weaker coherent organisation and more fragmented spatial structure during transition. By mode 40, the cumulative energy reaches approximately 54\% in \gls{CPG} and 83\% in \gls{TEQ}, demonstrating that \gls{TEQ} concentrates more of the total pressure-field energy into large-scale modal structures, whilst \gls{CPG} retains significant energy in unresolved small-scale features.

\begin{figure}[t]
    \centering
    \includegraphics[width=9.3cm]{resources/figures/breakdown/result_dns_thermal_pod_energy.pdf}
    \caption[Q-criterion isosurfaces coloured by density during transition for CPG, TEQ, and TNE models.]%
    {Fourier transform of pressure wall data.}
    \label{figure:pod_energy_variation}
\end{figure}


The Proper Orthogonal Decomposition (POD) analysis decomposes the pressure fluctuations into a set of spatial modes that capture the dominant coherent structures in the flow. The methodology begins by assembling the pressure field data from multiple timesteps into a spatiotemporal matrix of dimensions $(n_x \cdot n_z) \times n_t$, where $n_x$ and $n_z$ are the streamwise and spanwise grid points, and $n_t$ is the number of temporal snapshots. The data is mean-centred by subtracting the temporal average from each spatial location to isolate fluctuations. The Snapshot POD approach employs Singular Value Decomposition (SVD) on this centred matrix, which yields left singular vectors (spatial modes), right singular vectors (temporal coefficients), and singular values. The energy contained in each mode is quantified as $E_i = \lambda_i^2 / \sum \lambda_j^2$, where $\lambda_i$ represents the $i$-th singular value, allowing identification of the modes that capture the most energetic structures. The cumulative energy $CE_n = \sum_{i=1}^{n} E_i$ provides insight into the number of modes required to represent a given fraction of the total energy in the flow field. This decomposition facilitates direct comparison between the CPG and TEQ thermochemical models by examining whether the dominant modes and energy distribution differ significantly between the two approaches, thereby revealing differences in the coherent structure organisation across thermochemical assumptions.




The structure of POD mode 1, the most energetic single mode, represents the dominant coherent structure in the flow. Spatial analysis of mode 1 reveals quasi-periodic spanwise patterns that correspond directly to the streak structures identified in instantaneous and time-averaged Stanton-number distributions. The mode exhibits strong amplitude in the breakdown region ($x/\delta^{*}_{0} \approx 350$--$450$) and persists downstream into the turbulent region, reflecting the sustained spanwise organisation of heat transfer during transition. The spanwise wavelength captured by mode 1 is approximately $\lambda_z \approx 100$--$120\delta^{*}_{0}$, consistent with the subharmonic wavelengths introduced in the initial forcing and their subsequent nonlinear evolution. The mode 1 amplitude in \gls{TEQ} exceeds that in \gls{CPG}, indicating stronger coherent streak formation driven by the enhanced secondary instability growth in the thermally equilibrated case. In contrast, higher-order modes (e.g., mode 34) capture the turbulent fluctuations that develop in the fully turbulent region and the small-scale pressure variations associated with vortical structures. Mode 34 exhibits rapid spatial variation and lacks the organised periodic structure of mode 1, instead showing the signature of fine-scale turbulent eddies. The energies of mode 34 are comparable between \gls{CPG} and \gls{TEQ}, suggesting that once fully turbulent, the cascade of energy toward small scales proceeds at similar rates regardless of thermal modelling.

\begin{figure}[t]
    \centering
    \includegraphics[width=18.0cm,angle=-90]{resources/figures/breakdown/result_dns_thermal_contour_pod_mode1.pdf}
    \caption[Instantaneous Temperatsure ($T/T_e$) during boundary layer transition for CPG, TEQ, and TNE models.]%
    {Instantaneous streamwise velocity ($T/T_e$) during boundary layer transition for (a) \gls{CPG}, (b) \gls{TEQ}, and (c) \gls{TNE} models.}
    \label{figure:contour_podmode1}
\end{figure}


The progression of POD mode energies from mode 1 (large-scale streaks) through the intermediate modes to mode 34 (small-scale turbulence) illustrates the energy cascade mechanism during transition. The steeper energy decay observed in \gls{TEQ} (modes 1--10) compared to \gls{CPG} reflects the more coherent structure and concentration of energy in the large-scale modes. The flatter decay in \gls{CPG} (modes 1--20) indicates that energy is more uniformly distributed across scales, consistent with a less organised transition process. These findings confirm that thermal equilibrium enhances coherent structure formation, concentrating pressure fluctuations into energetic streak modes, whilst perfect-gas dynamics support a more distributed modal response. The POD analysis therefore provides a complementary perspective to the instantaneous isosurface observations and time-averaged Stanton-number distributions, directly linking the spectral properties of wall pressure to the spatial organisation and intensity of coherent structures during breakdown.




% \subsection{Boundary layer evolution}

% This section is where I need to include things about Boundary layer profiles and temperature profiles to show how the steep temperatures are effects the total features, maybe quantify here that the thermal nonequilbrium is pushed? or should I do this in the ILES chapter? 

\subsection{Boundary layer evolution}

The development of the boundary layer during transition provides insight into the changes in flow structure and momentum redistribution as instability modes grow and breakdown occurs. Integral parameters such as displacement thickness and momentum thickness characterise the overall profile development and the degree of flow blockage and momentum deficit. Examining these quantities across the three thermal models reveals how thermal modelling assumptions affect the rate of boundary layer thickening and the spatial distribution of transition stages.

The streamwise evolution of displacement thickness $\delta^*$ and momentum thickness $\theta$ is shown in Figure~\ref{figure:blgrowth}. In the laminar region ($x/\delta^*_0 < 200$), all three models exhibit similar growth rates, with $\delta^*$ and $\theta$ increasing gradually in accordance with Blasius-type scaling expected for a developing laminar boundary layer. The agreement across models in this region reflects the minimal influence of thermal effects when disturbance amplitudes remain small. As the second-mode resonance develops around $x/\delta^*_0 \approx 250$--$350$, the growth rates begin to diverge. \gls{TEQ} exhibits accelerated thickening consistent with its early transition onset, whilst \gls{CPG} shows intermediate growth. The \gls{TNE} case displays the most gradual thickening in this region, reflecting the delayed breakdown and attenuated instability growth rates predicted by linear studies. By the end of transition ($x/\delta^*_0 > 500$), all three models show accelerated growth as the flow becomes increasingly turbulent. The final displacement and momentum thicknesses converge toward comparable values, though the streamwise locations at which saturation is achieved differ across models, directly reflecting the transition delays observed in wall-bounded quantities. The shape factor $H = \delta^* / \theta$ provides additional insight into profile evolution, transitioning from approximately $13$--$14$ in the laminar state to below $11$ in fully developed turbulence, indicating progressive steepening of the velocity profile as coherent structures develop and breakdown proceeds.

\newpage
\begin{figure}[H]
    \centering
    \includegraphics[width=20.0cm,angle=-90]{resources/figures/breakdown/result_breadkdown_thermal_Tslices.pdf}
    \caption[Instantaneous Temperatsure ($T/T_e$) during boundary layer transition for CPG, TEQ, and TNE models.]%
    {Instantaneous streamwise velocity ($T/T_e$) during boundary layer transition for (a) \gls{CPG}, (b) \gls{TEQ}, and (c) \gls{TNE} models.}
    \label{figure:velocity_snapshot}
\end{figure}
\newpage


\subsection{Turbulent regime}
Once the boundary layer has undergone transition, the flow enters a fully turbulent regime characterised by intensified momentum and energy transport. In this stage, wall quantities such as skin friction and heat flux settle towards their asymptotic turbulent levels, while differences between thermal models primarily influence the turbulence structure and mean-flow development. To examine these effects, velocity statistics and Reynolds stresses are investigated. Due to the relatively small range of Reynolds numbers attained in the turbulent portion of the domain, the nondimensional flow profiles do not vary substantially along the flat plate. Therefore, the analysis is restricted to a single representative station. The corresponding streamwise locations for each model are listed in Table~\ref{table:global_quantities_models} in the turbulent region column.



\subsubsection{Velocity and temperature statistics}
A detailed assessment of velocity statistics is essential for characterising the turbulent boundary layer. Mean velocity profiles provide direct information on the redistribution of momentum by turbulence and form the basis for testing scaling laws and transformation approaches. Higher-order statistics, such as Reynolds stresses, quantify the intensity and anisotropy of turbulent fluctuations and reveal how energy is transferred among flow components. Together, these measures offer a comprehensive view of the turbulent structure and allow differences between thermal models to be identified and interpreted. In compressible boundary layers, direct comparison with incompressible velocity profiles requires a transformation that accounts for density and viscosity variations near the wall. The van Driest II transformation provides such a framework by extending the original van Driest scaling to include viscous effects, offering improved collapse of velocity profiles in the inner region. This makes it a common reference for assessing compressible turbulent boundary layers and for comparing different thermal models on a consistent basis.  The mean velocity and wall-normal coordinate $y$ are scaled using the friction velocity $u_{\tau}$ and the viscous length scale $\nu_w/u_{\tau}$. The van Driest transformation is commonly employed to compare compressible turbulent boundary layers with their incompressible counterparts. The transformed velocity is defined as  
\begin{equation}
    u^{+} = \int_{0}^{\bar{u}} \sqrt{\frac{\bar{\rho}}{\rho_w}} \, du ,
\end{equation}
where $\bar{\rho}$ is the time averaged local mean density and $\rho_w$ the density at the wall.  

The resulting mean velocity profile should collapse onto the universal law of the wall, matching the viscous sublayer ($u^{+}=y^{+}$ for $y^{+}<10$) and approaching the logarithmic law at larger wall-normal distances:
\begin{equation}
    u^{+} = \frac{1}{\kappa}\ln(y^{+}) + C, 
    \qquad \kappa = 0.41, \quad C = 5.2 .
\end{equation}

Recent work has introduced updated velocity transformations that improve the collapse of compressible profiles, particularly under cold-wall conditions~\cite{HasanEtal2023}. These approaches build on the semilocal scaling ideas of Trettel and Larsson~\cite{TrettelLarsson2016}, which emphasise the role of variable density and viscosity in defining an effective viscous length scale.

To define the transformation, a consistent relationship between $Y^{+}$ and $y^{+}$ is required. Under the assumption of universality of the turbulent shear stress in the near-wall region, the relation
\begin{equation}
    Y^{+} = \frac{\delta_{v}}{\delta^{*}_{v}}\, y^{+} = y^{*}
\end{equation}
is introduced. In practice, this universality is not guaranteed once compressibility modifies the near-wall stress distribution, so the equivalence $Y^{+}=y^{*}$ requires reassessment. The interpretation
\begin{equation}
    Y^{+} = \frac{\nu}{\delta^{*}_{v}}
\end{equation}
identifies $\delta^{*}_{v}$ as the relevant viscous length scale for small-scale motions in compressible flows and forms the basis of the HLPP formulation.

Making use of the coordinate mapping $Y^{+}=y^{*}$, together with
\begin{equation}
    \frac{dy^{*}}{dy} = \left(1-y^{*}\frac{d\delta^{*}_{v}}{dy}\right)\frac{\delta_{v}}{\delta^{*}_{v}},
    \qquad
    \frac{u_{\tau}}{u^{*}_{\tau}}=\sqrt{\frac{\bar{\rho}}{\rho_{w}}},
\end{equation}
the transformation kernel becomes
\begin{equation}
\frac{d\overline{U}^{+}}{dy^{+}} = 
\left(\frac{1+\mu^{c}_{t}/\bar{\mu}}{1+\mu^{i}_{t}/\mu_{w}}\right)^{1/3}
\left(1-\frac{y}{\delta^{*}_{v}}\frac{d\delta^{*}_{v}}{dy}\right)^{1/2}
\sqrt{\frac{\bar{\rho}}{\rho_{w}}},
\label{eq:TLkernel}
\end{equation}
where the three multiplicative terms represent adjustments to (i) friction velocity scaling, (ii) the viscous length scale, and (iii) density variations in the inner layer.

To close the formulation, explicit models for the eddy viscosities $\mu^{i}_{t}$ and $\mu^{c}_{t}$ are introduced. The Johnson–King model with Van Driest-type damping is adopted, giving
\begin{equation}
    \mu^{i}_{t}=\sqrt{\tau_{w}\rho_{w}}\,\kappa y D^{i},
    \qquad
    D^{i}=\left[1-\exp\left(-\frac{y^{+}}{A^{+}}\right)\right]^{2},
\end{equation}
with $\kappa=0.41$ and $A^{+}=17$ yielding a log-law intercept close to 5.2. For compressible flows, the corresponding form is
\begin{equation}
    \mu^{c}_{t}=\sqrt{\tau_{w}\rho_{w}}\,\kappa y D^{c},
\end{equation}
where $D^{c}$ is a damping function based on the semilocal coordinate $y^{*}$.

Integrating the kernel gives the transformed velocity profile,
\begin{equation}
u^{+}_{HLPP}(y^{*}) = \int_{0}^{y^{*}}
\left(\frac{1+\mu^{c}_{t}/\bar{\mu}}{1+\mu^{i}_{t}/\mu_{w}}\right)^{1/3}
\left(1-\frac{y}{\delta^{*}_{v}}\frac{d\delta^{*}_{v}}{dy}\right)^{1/2}
\sqrt{\frac{\bar{\rho}}{\rho_{w}}}\, dy^{*},
\label{eq:hlpp}
\end{equation}
where $u^{+}_{HLPP}$ denotes the transformed velocity in wall units. This representation allows mean velocity profiles in compressible boundary layers to be compared consistently across thermal models and remains accurate even under cold-wall conditions.


\begin{figure}[t]
    \centering
    \fontsize{9}{10}\selectfont
    \includegraphics[width=1.04\textwidth]{resources/figures/breakdown/result_thermal_velocity_temperature_statistics.pdf}
    \caption[Velocity and temperature statistics in transformed wall units for Mach~6, $T_w=1800$ K.]%
    {Profiles of velocity and temperature statistics plotted against the semilocal wall coordinate $y^{*}$, obtained from the van Driest transformation, for the Mach~6, $T_w=1800$ K case. Results are shown for \gls{CPG}, \gls{TEQ}, and \gls{TNE} formulations.}
    \label{figure:thermal_statistics}
\end{figure}

As shown in Figure~\ref{figure:thermal_spanwise_st}, the velocity profiles collapse well onto the expected scaling laws when expressed in terms of the semilocal wall coordinate $y^{*}$. The curves follow the linear relation in the viscous sublayer and approach the logarithmic law in the overlap region, confirming that the HLPP transformation provides a consistent framework across the different thermal models.  

The temperature profiles, however, do not show full equilibration. Whereas linear stability theory indicated that in the downstream region the \gls{TNE} mean temperature tends to align more closely with \gls{TEQ} than with \gls{CPG}, the present turbulent results reveal the opposite behaviour: turbulence slows down the equilibration process, and the \gls{TNE} profiles remain distinctly separated from \gls{TEQ} further into the boundary layer. This highlights the role of turbulence in delaying thermal relaxation compared to the laminar or transitional regimes. Overall, these results demonstrate that while velocity statistics conform to classical scaling once turbulence is established, thermal nonequilibrium effects persist much further downstream, underscoring the fundamentally different relaxation behaviour of velocity and temperature fields in turbulent boundary layers.


\begin{figure}[t]
    \centering
    \fontsize{9}{10}\selectfont
    \includegraphics[width=1.04\textwidth]{resources/figures/breakdown/result_thermal_reynolds_stresses.pdf}
    \caption[Reynolds stress components in semilocal wall units for Mach~6, $T_w=1800$~K.]%
    {Reynolds stress components $\overline{\rho}\,\widetilde{u''u''}$, 
     $\overline{\rho}\,\widetilde{v''v''}$, 
     $\overline{\rho}\,\widetilde{w''w''}$, and 
     $-\overline{\rho}\,\widetilde{u''v''}$ plotted against the semilocal wall coordinate $y^{*}$ 
     for the Mach~6, $T_w=1800$~K case. Results are compared between \gls{CPG}, \gls{TEQ}, and \gls{TNE} formulations at multiple streamwise stations, with the final location highlighted.}
    \label{figure:thermal_reynolds_stresses}
\end{figure}

\subsubsection{Reynolds stresses}

The Reynolds stress components at $x/\delta^{*}_{0}=650$ and $700$, extracted from Favre-averaged statistics and normalised by the mean wall shear stress $\tau_{w}$, exhibit the canonical ordering $\overline{\rho}\,\widetilde{u''u''} > \overline{\rho}\,\widetilde{v''v''} \sim \overline{\rho}\,\widetilde{w''w''}$, with the shear stress $-\overline{\rho}\,\widetilde{u''v''}$ peaking closer to the wall, as shown in figure~\ref{figure:thermal_reynolds_stresses}. The streamwise component reaches $\overline{\rho}\,\widetilde{u''u''}/\tau_{w}\approx 14$ at $y^{*}\approx 12$, while $\overline{\rho}\,\widetilde{v''v''}$ and $\overline{\rho}\,\widetilde{w''w''}$ attain peak values of $\mathcal{O}(1.0$–$1.5)$, and the shear stress reaches a maximum of about unity. These observations are consistent with canonical DNS of hypersonic turbulent boundary layers \citep{ZhangDuanChoudhari2018,XuEtal2021}. At both stations, the \gls{CPG}, \gls{TEQ}, and \gls{TNE} closures collapse onto nearly indistinguishable profiles, demonstrating that once turbulence is fully developed, the inclusion of thermal nonequilibrium effects does not significantly alter the self-similar Reynolds stress structure.

\begin{figure}[t]
    \centering
    \fontsize{9}{10}\selectfont
    \includegraphics[width=1.04\textwidth]{resources/figures/breakdown/result_thermal_reynolds_rho_T.pdf}
    \caption[Density and temperature statistics in semilocal wall units for Mach~6, $T_w=1800$~K.]%
    {Profiles of Favre-averaged density (top row) and temperature (bottom row) together with their corresponding normalised root-mean-square fluctuations, plotted against the semilocal wall coordinate $y^{*}$ for the Mach~6, $T_w=1800$~K case. Blue lines correspond to the \gls{CPG} formulation and red lines to the \gls{TNE} formulation, with results shown at several streamwise stations and the final location emphasised.}
    \label{figure:thermal_rho_T}
\end{figure}

Figure~\ref{figure:thermal_rho_T} presents the evolution of density and temperature statistics across the boundary layer. In the mean profiles (left panels), both \gls{CPG} and \gls{TNE} display broadly similar behaviour, with the density and temperature ratios decreasing from elevated near-wall values towards the freestream. Subtle differences emerge downstream, where the \gls{TNE} solutions relax more gradually, reflecting the delayed equilibration associated with vibrational energy modes.

More pronounced deviations are observed in the fluctuation levels (right panels). Within the turbulent region, the \gls{TNE} formulation consistently exhibits larger root-mean-square variations in both density and temperature compared to the \gls{CPG} case. This amplification arises from the additional degrees of freedom introduced by vibrational nonequilibrium, which enhance the coupling between thermal and momentum fluctuations. The resulting increase in density and temperature variances strengthens turbulent mixing and alters the redistribution of energy within the boundary layer. In the mean temperature profiles, the \gls{CPG} and \gls{TNE} cases exhibit similar near-wall behaviour, but the \gls{TNE} solution decays more gradually towards the freestream, reflecting delayed vibrational relaxation. These findings indicate that, even when mean profiles appear to converge, the turbulence structure remains sensitive to the inclusion of vibrational nonequilibrium, with direct consequences for predicting wall heat transfer and informing model development.



% \section{Grid Study of DNS and ILES}
% In the following section, there will be a grid study of ILES cases to DNS cases, 

% \begin{figure}[t]
%     \centering
%     \fontsize{9}{10}\selectfont
%     \includegraphics[width=0.85\textwidth]{resources/figures/test/result_wallquantities_cf_st.pdf}
%     \caption[Streamwise evolution of $C_f$ and $St$ for Mach~6, $T_w=1800$ K.]%
%     {Streamwise evolution of spatially temporally averaged $C_f$ and $St$ for the Mach~6, $T_w=1800$ K case. Results are shown for \gls{CPG}, \gls{TEQ}, and \gls{TNE} formulations, alongside laminar reference solutions and the Van Driest~II turbulent correlation.}
%     \label{figure:gridstudy_cfst}
% \end{figure}

% % \subsection{Transitional regime}

% For that reason, the thesis concludes that ILES cases, although not good at detecting the turbulence intensity, are good at getting the transition onset location prediction and the nonlinear growth of the 2D modes. 

% \section{Role of Mach number and wall thermal conditions}

% \subsubsection{Results}
% \subsection{Breakdown to turbulence}
% Turbulent mean and flow fluctuations
% \subsubsection{Thermal nonequilibrium effects}
% \subsubsection{Heat flux and Reynolds stresses}

% \section{Role of Mach number and wall thermal conditions}

% % \subsection{Results}
% \subsection{Global flow properties}
% \subsection{Boundary-layer profiles}
% % \subsection{Shock--boundary layer interaction}

% \subsection{Breakdown to turbulence}
% \subsubsection{Onset location and growth rates}
% \subsubsection{Mode interactions and energy transfer}
% \subsubsection{Thermal nonequilibrium effects}
% \subsubsection{Heat flux and Reynolds stresses}
% \subsubsection{Second-order statistics}
% % \subsubsection{Spectral analysis of turbulence}

% \newpage

\section{Role of Mach number and wall thermal conditions}
The combined influence of Mach number and wall thermal conditions is now assessed within the framework of thermal nonequilibrium. These parameters exert a strong control on the stability and transition characteristics of hypersonic boundary layers: increasing Mach number enhances compressibility effects and amplifies acoustic-trapped disturbances, while variation in wall temperature modifies the mean thermal gradients that drive instability growth. To isolate these effects, simulations are performed for \gls{TNE} flows at Mach~6 and Mach~8, each with both a cold wall at $T_w=1200$~K and a hot wall at $T_w=1800$~K. This configuration enables a systematic examination of how freestream Mach number and wall thermal state interact to influence instability amplification, breakdown mechanisms, and the resulting surface heating environment.


\subsection{Global flow properties}

Here is where we have the skin friction plots, Stanton numbers and then the big table that no one asked about:

\begin{figure}[h]
    \centering
    \fontsize{9}{10}\selectfont
    \includegraphics[width=0.95\textwidth]{resources/figures/breakdown/result_wallquantities_cf_tne.pdf}
    \caption[Streamwise evolution of $C_f$ and $St$ for Mach~6, $T_w=1800$ K.]%
    {Streamwise evolution of spatially temporally averaged $C_f$ and $St$ for the Mach~6, $T_w=1800$ K case. Results are shown for \gls{CPG}, \gls{TEQ}, and \gls{TNE} formulations, alongside laminar reference solutions and the Van Driest~II turbulent correlation.}
    \label{figure:gridstudy_cfst}
\end{figure}

\begin{figure}[h]
    \centering
    \fontsize{9}{10}\selectfont
    \includegraphics[width=0.95\textwidth]{resources/figures/breakdown/result_wallquantities_st_tne.pdf}
    \caption[Streamwise evolution of $C_f$ and $St$ for Mach~6, $T_w=1800$ K.]%
    {Streamwise evolution of spatially temporally averaged $C_f$ and $St$ for the Mach~6, $T_w=1800$ K case. Results are shown for \gls{CPG}, \gls{TEQ}, and \gls{TNE} formulations, alongside laminar reference solutions and the Van Driest~II turbulent correlation.}
    \label{figure:gridstudy_cfst}
\end{figure}




\subsection{Breakdown to turbulence}

Understanding the breakdown process, which is significantly delayed in the Mach 8 cases, what causes this breakdown? 

\subsection{Onset location and growth rates}
% \subsubsection{Boundary-layer profiles}

\subsection{Mode interactions and energy transfer}

Similar work of POD from last chapter, but with some extra results to show dominant modes of Mack 8 and hopefully answer the question of why there is a dip in Mach 6 cases and not the Mach 8 case. 

\subsection{Streak breakdown}
The main section to talk about I guess, Streak amplitude plot, breakdown to turbulence due to the prescence of streaks, \citep{Andersson2001} and \citep{Passiatore2021} have spoke  about this. 


\subsection{Thermal nonequilibrium effects}
This is where we introduce the thermal nonequilibrium effects from statistics to how the delay in transition happens, 

\subsubsection{Heat flux and Reynolds stresses}
Similar story of ILES domain getting this correct

\subsubsection{Second-order statistics}


\section{Chapter Summary}


